\section{Introduction}
\label{sec:intro}

A Zero Trust architecture (ZTA) replaces the perimeter with per-request
authorization: every access is evaluated against identity, device posture and
resource sensitivity, continuously, for the lifetime of a session~\cite{nist800207}.
This shifts the detection problem. The question is no longer whether a packet crossed
a boundary, but whether \emph{this} principal, arriving from \emph{this} client, on
\emph{this} device, has any history of touching \emph{that} asset.

Lateral movement is the class for which that question is hardest and matters most. An
attacker who has obtained valid credentials generates traffic that is, feature by
feature, indistinguishable from legitimate traffic: the request is authenticated, the
policy permits it, and no sensor fires. What separates it from benign activity is
relational and temporal: the pair has no history, and the access follows a
reconnaissance burst on the same origin. Detecting it as anomalous-edge detection over
a temporal graph is an established formulation, with \textit{Euler}~\cite{euler} and \textit{Argus}~\cite{argus} as reference systems and the LANL authentication
corpus~\cite{lanl} as the benchmark. While formulating lateral movement as dynamic
graph link prediction is well explored, existing systems operate on substrates
that discard rich contextual signals that a ZTA policy decision point (PDP) already possesses.

The gap is concrete. A host-to-host authentication log records that computer $A$
authenticated to computer $B$ as user $u$. A ZTA request carries, in addition, the
client network origin, the TLS fingerprint of the issuing client application, the
attested device identity, and the specific asset requested. Collapsing those into
$(A, B, u)$ makes 1 attack class \emph{inexpressible}: an adversary who replays a
stolen credential from an external machine produces a request whose principal is
familiar but whose client configuration is novel. On an authentication graph there is
no node to capture that discrepancy.

\subsection*{Contributions}

\begin{enumerate}
  \item \textit{A 5-node causal chain per request.} Each ZTA access is decomposed
  into \textit{source} $\rightarrow$ \textit{config} $\rightarrow$ \textit{device}
  $\rightarrow$ \textit{user} $\rightarrow$ \textit{resource}, with an explicit
  \textit{config}\,$\rightarrow$\,\textit{user} binding edge. The \textit{config}
  node, keyed by the client TLS fingerprint (JA3), enables detection of credential reuse
  from an unfamiliar tool as a structural anomaly (Section~\ref{sec:approach}).

  \item \textit{An external policy-driven commit gate.} To mitigate self-poisoning
  in streaming settings, entity memory and temporal neighbourhoods advance only on
  events admitted by an external policy engine, decoupling memory retention from the
  detector's own predictions (Section~\ref{sec:problem}).

  \item \textit{A rigorous leakage audit for synthetic benchmarks.} To prevent
  unintended feature shortcuts common in synthetic benchmark generation, a regression
  audit is developed that bounds single-feature separability, forbids exact-value
  fingerprints, and equalizes destination marginals across benign and anomalous classes
  (Section~\ref{sec:setup}).

  \item \textit{An external-validity study on real enterprise telemetry.} A survey of
  public cybersecurity datasets identifies the absence of corpora uniting user identity
  and TLS fingerprints. On PicoDomain~\cite{picodomain}, the only corpus providing both,
  the schema adapter is implemented, session-binding coverage is quantified, and
  transferability is evaluated (Section~\ref{sec:external}).
\end{enumerate}

Section~\ref{sec:related} positions the work against the temporal-graph lateral
movement literature. Section~\ref{sec:problem} states the threat model and the
learning problem, Section~\ref{sec:approach} describes the architecture,
Section~\ref{sec:setup} details the experimental protocol, Section~\ref{sec:results} presents
the empirical evaluation, and Section~\ref{sec:limits} discusses known limitations.
